% =============================================================
% บทที่ 3 วิธีการดำเนินงาน
% =============================================================

% -----------------------------------------------------------
\section{ขั้นตอนการดำเนินงาน}
% -----------------------------------------------------------

การดำเนินโครงงานระบบรวบรวมและสังเคราะห์ประเด็นข่าวภาษาไทยอัตโนมัติ (End-to-End News Collector and Summarizer) ใช้ระเบียบวิธีวิศวกรรมซอฟต์แวร์เชิงขั้นตอนแบบวนซ้ำ (Iterative Software Engineering Methodology) เพื่อให้การวิเคราะห์ ออกแบบ พัฒนา และทดสอบระบบเป็นไปอย่างมีแบบแผนและสามารถตรวจสอบย้อนกลับได้ทุกขั้นตอน โดยแบ่งกระบวนการดำเนินงานออกเป็น 5 ขั้นตอนหลัก ดังนี้

\subsection{ศึกษาปัญหาและรวบรวมความต้องการ}

คณะผู้จัดทำได้ศึกษาสภาพปัญหาและพฤติกรรมการติดตามข่าวสารออนไลน์ของผู้ใช้ในยุคดิจิทัลผ่าน 2 กระบวนการหลัก ได้แก่

\begin{enumerate}
    \item \textbf{การสัมภาษณ์แบบกึ่งโครงสร้าง (Semi-structured Interview)}\\*
    ดำเนินการสัมภาษณ์กลุ่มผู้ติดตามข่าวสารออนไลน์ เป็นนักศึกษา 5 คน เพื่อรวบรวมพฤติกรรมการเสพข่าว และความต้องการด้านฟีเจอร์ที่พึงประสงค์
    \item \textbf{การสังเกตการณ์และการสำรวจระบบงานเดิม (Observation and Baseline Survey)}\\*
    สำรวจการนำเสนอข่าวและโครงสร้างของสำนักข่าวออนไลน์ภาษาไทยชั้นนำ เช่น ไทยรัฐ, มติชน, ข่าวสด, ผู้จัดการออนไลน์, ประชาไท, The Standard, เดลินิวส์, สำนักข่าวอิศรา, PPTV HD 36, Thai PBS, Workpoint Today, กรุงเทพธุรกิจ, โพสต์ทูเดย์ และ Sanook News เพื่อศึกษารูปแบบ RSS Feed โครงสร้าง HTML และการใช้งานเทคโนโลยีเรนเดอร์หน้าเว็บ
\end{enumerate}

จากการศึกษาและรวบรวมข้อมูล คณะผู้จัดทำพบปัญหาสำคัญ 2 ประการหลัก ได้แก่

ประการแรก ปัญหาความกระจัดกระจายและข้อมูลข่าวสารซ้ำซ้อน (Information Overload and Fragmentation) เนื่องจากสำนักข่าวแต่ละแห่งต่างรายงานข่าวเดียวกันด้วยมุมมองและสำนวนพาดหัวที่แตกต่างกัน ทำให้ผู้อ่านต้องเปิดหลายเว็บไซต์และเสียเวลาคัดกรองเนื้อหาซ้ำซ้อน

ประการสอง ปัญหาผู้ใช้มักชอบอ่านแค่พาดหัวข่าว ทำให้ขาดความเข้าใจประเด็นอื่นๆ ในเนื้อหาข่าวนั้นๆ

ประการที่สาม ปัญหาความไม่ต่อเนื่องของประเด็นข่าว (Loss of Story Context) ข่าวเหตุการณ์ใหญ่ที่มีความยืดเยื้อข้ามสัปดาห์ เช่น อุทกภัยหรือคดีความ มักถูกตัดตอนและแยกส่วนตามกรอบเวลา 24--48 ชั่วโมง ทำให้ผู้อ่านขาดความเข้าใจในพัฒนาการตามลำดับเวลา

\subsection{วิเคราะห์และออกแบบระบบ}

คณะผู้จัดทำได้นำผลการศึกษาปัญหาและความต้องการมากำหนดเป็นข้อกำหนดระบบที่วัดผลได้อย่างเป็นรูปธรรม ประกอบด้วย

\begin{enumerate}
    \item \textbf{การกำหนดความต้องการเชิงฟังก์ชัน (Functional Requirements)}\\*
    กำหนดความสามารถของระบบครอบคลุม 15 ข้อกำหนดหลัก (FR-01 ถึง FR-15) ได้แก่ การรวบรวมข่าวสารจากหลายสำนัก การลงทะเบียนแหล่งข่าวที่สามารถขยายต่อได้ การจำแนก 9 หมวดหมู่ภาษาไทย การสรุปเนื้อหาและวิเคราะห์อารมณ์ การนำเสนอแบบ Bento Hub การกรองและการค้นหา การพรีวิวเนื้อหาบทความ การบันทึกคั่นหน้า การระบุประเด็นข่าวสองชั้น การจัดอันดับกระแสข่าวสาร การเชื่อมโยงเส้นเรื่องมหาภาค แฟ้มสังเคราะห์ไทม์ไลน์และข้อเท็จจริงยืนยัน กลไกควบคุมความถูกต้อง 3 ระดับ ศูนย์ควบคุมสถานะของผู้ดูแลระบบ และการคุ้มครองความเป็นส่วนตัวตามมาตรฐาน PDPA
    \item \textbf{การกำหนดความต้องการที่ไม่ใช่เชิงฟังก์ชัน (Non-functional Requirements)}\\*
    กำหนดเกณฑ์คุณภาพทางวิศวกรรม 8 มิติ (NFR-01 ถึง NFR-08) ได้แก่ เวลาตอบสนองคำขอต่ำกว่า 150 มิลลิวินาที ความหน่วงการกระจายข้อมูลสดต่ำกว่า 500 มิลลิวินาที อัตราการดึงข้อมูลไม่น้อยกว่า 100 บทความต่อนาที การแยกความเสียหายข้ามสำนักข่าว ความพร้อมใช้งานไม่ต่ำกว่าร้อยละ 99.5 ความแม่นยำ B-cubed Precision $\ge 0.85$ และ B-cubed F1 $\ge 0.75$ การรักษาความปลอดภัยตามหลัก Zero-PII และการทำงานข้ามแพลตฟอร์ม
    \item \textbf{การกำหนดกฎการตรวจสอบความถูกต้องของข้อมูล (Validation Rules)}\\*
    กำหนดกฎเกณฑ์ควบคุมความสมบูรณ์ของข้อมูล 10 กฎ (VR-01 ถึง VR-10) ผ่านโมเดล Pydantic เพื่อป้องกันข้อผิดพลาดในการประมวลผลและการจัดเก็บ
    \item \textbf{การจัดทำแบบจำลองระบบ (System Models)}\\*
    จัดทำแผนภาพกรณีใช้งาน (Use Case Diagram) และข้อกำหนดกรณีใช้งาน (UC-01 ถึง UC-04), แผนภาพลำดับขั้นตอน (Sequence Diagram) 3 กระบวนการหลัก, สถาปัตยกรรมเชิงตรรกะ 4 ชั้น (Logical Architecture), แผนภาพการติดตั้งระบบ (Deployment Architecture), แผนภาพสถานะ (State Machine Diagrams) ของวงจรชีวิตประเด็นข่าวและโหมดการสังเคราะห์, แผนภาพความสัมพันธ์ของข้อมูล (Relational Diagram) สำหรับ 6 ตารางหลัก และต้นแบบส่วนติดต่อกับผู้ใช้ (UI Mockup) จำนวน 6 รูปแบบ (M-01 ถึง M-06)
\end{enumerate}

\subsection{พัฒนาระบบ}

การพัฒนาระบบแบ่งการทำงานออกเป็น 4 รอบการทำงานหลัก (Iterations) เพื่อให้สามารถทดสอบและปรับปรุงคุณภาพของแต่ละส่วนได้อย่างต่อเนื่อง ดังนี้

\begin{enumerate}
    \item \textbf{รอบการทำงานที่ 1 การรวบรวม สกัดเนื้อหา และจำแนกหมวดหมู่ข่าวสาร}\\*
    พัฒนาไปป์ไลน์การดึงข่าวสารอัตโนมัติแบบอะซิงโครนัสด้วย HTTPX และ Playwright Headless Browser, การสกัดเนื้อหาบทความด้วย Trafilatura, ระบบแคชหลายระดับ (Multi-tier Caching) สองชั้น และโมเดลจำแนกหมวดหมู่ข่าวภาษาไทย 9 หมวดหมู่อย่างรวดเร็วด้วยสถาปัตยกรรม 3-Tier Cascade Classifier (Regex Cues, LinearSVC/TF-IDF และ WangchanBERTa)
    \item \textbf{รอบการทำงานที่ 2 การระบุประเด็นข่าวออนไลน์ 2 ชั้นและการจัดอันดับกระแสข่าว}\\*
    พัฒนาระบบระบุประเด็นข่าวออนไลน์แบบคงทน (Online Persistent Story Identification) ด้วยเวกเตอร์ฝังตัว WangchanBERTa ผสานการตรวจสอบเอนทิตีชื่อเฉพาะ thainer-v2 ผ่านเกณฑ์ Two-Band Cosine Gate และปรับปรุง Centroid แบบ EMA ข้ามหน้าต่าง 72 ชั่วโมง พร้อมพัฒนาเอนจินจัดอันดับกระแสข่าวสาร 6 มิติเชิงสถิติแบบ Zero-LLM ที่ตัดลดทอนข่าวประจำวัน (Routine Penalty)
    \item \textbf{รอบการทำงานที่ 3 แฟ้มสังเคราะห์ความคืบหน้าจากหลักฐานจริงและระบบเจาะประเด็นข่าว}\\*
    พัฒนากระบวนการทำงานแบบ 3 ขั้นตอนที่บังคับให้ใช้ผลการคัดกรอง (Enforced Multi-Stage Pipeline) ประกอบด้วย Event Relevance Gate, Progression Timeline และ Corroborated Facts ที่ผ่านการตรวจสอบรหัสอ้างอิงรายประโยคอย่างเข้มงวด พร้อมพัฒนาระบบสืบค้นเจาะประเด็นข่าวเชิงลึก (Issue Explorer) บนฐานข้อมูล SQLite ในโหมดอ่านอย่างเดียว (Read-Only) และระบบปิดการทำงานฉุกเฉิน (Kill Switch)
    \item \textbf{รอบการทำงานที่ 4 ส่วนต่อประสานผู้ใช้ การสื่อสารเรียลไทม์ และระบบควบคุม}\\*
    พัฒนาหน้าจอ Bento Hub SPA ด้วย Vanilla JavaScript แบบโมดูลาร์ (Modular ES Modules) และ Tailwind CSS, การเชื่อมต่อผลักดันข้อมูลสดสองทิศทางผ่าน Socket.IO, หน้าต่างพรีวิวเนื้อหาข่าว In-App Reader Modal และการสังเคราะห์สรุปด้วย AI, ระบบบันทึกคั่นหน้าและประวัติการอ่านบน LocalStorage, ระบบจัดการความยินยอมคุกกี้ตามมาตรฐาน PDPA และศูนย์ควบคุมสถานะของผู้ดูแลระบบ (Ops Console)
\end{enumerate}

\subsection{ทดสอบและประเมินผล}

การทดสอบและประเมินผลระบบประกอบด้วยการทดสอบทางวิศวกรรมซอฟต์แวร์และการประเมินความพึงพอใจของผู้ใช้งานจริง โดยแบ่งออกเป็น 2 ส่วนหลัก ดังนี้

ประการแรก การทดสอบการยอมรับของผู้ใช้ (User Acceptance Testing หรือ UAT) ดำเนินการทดสอบตามกรณีทดสอบ 15 กรณี (TC-01 ถึง TC-15) ซึ่งครอบคลุมความต้องการเชิงฟังก์ชันครบทุกข้อ (FR-01 ถึง FR-15) โดยให้กลุ่มตัวอย่างผู้ใช้งานจริงดำเนินการทดสอบตามขั้นตอนอย่างเป็นระบบในภาคผนวก ข และคำนวณร้อยละของกรณีทดสอบที่ผ่านตามสมการที่~\eqref{eq_uat_percent} ซึ่งอ้างอิงตามเกณฑ์การประเมินการทดสอบซอฟต์แวร์มาตรฐานของ Pressman และ Maxim~\cite{pressman2014} ร่วมกับแนวทางของ Sommerville~\cite{sommerville2016}
 
\begin{equation}
\label{eq_uat_percent}
P = \frac{n_p}{N} \times 100
\end{equation}
เมื่อ
\begin{itemize}[noitemsep,topsep=0pt]
    \item $P$ คือ ร้อยละของกรณีทดสอบที่ผ่าน
    \item $n_p$ คือ จำนวนกรณีทดสอบที่ผ่านการทดสอบ
    \item $N$ คือ จำนวนกรณีทดสอบทั้งหมด ($N = 15$)
\end{itemize}

ประการที่สอง การประเมินความพึงพอใจของผู้ใช้งาน ดำเนินการเก็บรวบรวมข้อมูลจากกลุ่มตัวอย่างผู้ใช้งานจริงที่ได้ทดลองใช้งานระบบผ่านเว็บไซต์จริง (https://e2enews.wittensky.com) จำนวน 9 คน ผ่านแบบสอบถามประเมินความพึงพอใจและประสิทธิภาพการใช้งานระบบบนระบบออนไลน์ Google Forms ตามภาคผนวก ค โดยแบบสอบถามแบ่งออกเป็น 3 ส่วนหลัก ได้แก่ ส่วนที่ 1 ข้อมูลทั่วไปและพฤติกรรมการติดตามข่าวสารของผู้ตอบแบบสอบถาม, ส่วนที่ 2 การประเมินฟังก์ชันการทำงานหลัก (Core Features Evaluation) ครอบคลุม 6 ด้านสำคัญผ่านเกณฑ์ Rating Scale 5 ระดับ (1 หมายถึง ไม่พึงพอใจเลย ถึง 5 หมายถึง พึงพอใจมากที่สุด) และส่วนที่ 3 ความพึงพอใจโดยรวม ประสบการณ์การใช้งาน และการบอกต่อ (Overall Satisfaction and Net Promoter Score หรือ NPS) พร้อมข้อคำถามปลายเปิดเพื่อรับฟังปัญหาและข้อเสนอแนะในการพัฒนา โดยมีเกณฑ์การแปลความหมายค่าเฉลี่ยตามตารางที่~\ref{tab_likert_criteria}

\begin{table}[H]
\centering
\caption{เกณฑ์การแปลความหมายค่าเฉลี่ยความพึงพอใจ}
\label{tab_likert_criteria}
\begin{tabular}{|C{3.5cm}|C{4.5cm}|}
\hline
\textbf{ช่วงค่าเฉลี่ย} & \textbf{ระดับความพึงพอใจ} \\
\hline
4.51--5.00 & ดีมาก \\ \hline
3.51--4.50 & ดี \\ \hline
2.51--3.50 & ปานกลาง \\ \hline
1.51--2.50 & พอใช้ \\ \hline
1.00--1.50 & ควรปรับปรุง \\
\hline
\end{tabular}
\end{table}

การวิเคราะห์ผลคะแนนความพึงพอใจใช้การคำนวณค่าเฉลี่ยเลขคณิต ($\bar{x}$) ตามสมการที่~\eqref{eq_mean} และค่าส่วนเบี่ยงเบนมาตรฐาน ($\text{S.D.}$) ตามสมการที่~\eqref{eq_std} ซึ่งอ้างอิงตามหลักการวัดและประเมินผลทางการวิจัยของ Best และ Kahn~\cite{best2006}

\begin{equation}
\label{eq_mean}
\bar{x} = \frac{1}{n} \sum_{i=1}^n x_i
\end{equation}

\begin{equation}
\label{eq_std}
\text{S.D.} = \sqrt{\frac{\sum_{i=1}^n (x_i - \bar{x})^2}{n - 1}}
\end{equation}
เมื่อ
\begin{itemize}[noitemsep,topsep=0pt]
    \item $\bar{x}$ คือ ค่าเฉลี่ยของคะแนนความพึงพอใจ
    \item $\text{S.D.}$ คือ ส่วนเบี่ยงเบนมาตรฐานของคะแนนความพึงพอใจ
    \item $x_i$ คือ คะแนนความพึงพอใจของผู้ตอบแบบสอบถามคนที่ $i$ (ช่วงคะแนน 1--5)
    \item $n$ คือ จำนวนผู้ตอบแบบสอบถามทั้งหมด ($n = 9$)
\end{itemize}

\subsection{ปรับปรุงและสรุปผล}

นำผลการทดสอบการยอมรับของผู้ใช้และข้อเสนอแนะจากการประเมินความพึงพอใจมาปรับปรุงระบบ ได้แก่ การปรับแต่งเกณฑ์ค่าความคล้ายคลึงของประเด็นข่าว การขยายขีดความสามารถของระบบแคชเพื่อลดต้นทุนการเรียกใช้โมเดลภาษาขนาดใหญ่ การปรับแต่งการแสดงผลส่วนติดต่อผู้ใช้ให้รองรับอุปกรณ์เคลื่อนที่ และการจัดทำเอกสารคู่มือการติดตั้งตลอดจนรายงานโครงงานฉบับสมบูรณ์

% -----------------------------------------------------------
\section{เครื่องมือที่ใช้ในการดำเนินงาน}
% -----------------------------------------------------------

การดำเนินโครงงานจำเป็นต้องใช้เครื่องมือทั้งด้านฮาร์ดแวร์และซอฟต์แวร์ที่มีประสิทธิภาพ เพื่อรองรับการประมวลผลข้อมูลภาษาธรรมชาติและการจำลองเบราว์เซอร์ โดยมีรายละเอียดดังนี้

\subsection{ฮาร์ดแวร์}

เครื่องมือด้านฮาร์ดแวร์ที่ใช้ในการพัฒนาและทดสอบระบบแสดงดังตารางที่~\ref{tab_hardware}

\begin{table}[H]
\centering
\caption{ฮาร์ดแวร์ที่ใช้ในการพัฒนาและทดสอบ}
\label{tab_hardware}
\resizebox{\textwidth}{!}{%
\begin{tabular}{|c|L{4.5cm}|L{8.5cm}|}
\hline
\textbf{ลำดับ} & \textbf{รายการ} & \textbf{รายละเอียดคุณลักษณะ} \\
\hline
1 & เครื่องคอมพิวเตอร์สำหรับการพัฒนา & หน่วยประมวลผล Intel Core Ultra 5, หน่วยความจำหลัก RAM 16 GB, หน่วยบันทึกข้อมูล SSD ขนาด 512 GB \\ \hline
2 & เครื่องแม่ข่ายสำหรับทดสอบบนคลาวด์ & AWS EC2 Instance (c7i-flex.large), Ubuntu 22.04 LTS, ขนาด 2 vCPU, หน่วยความจำหลัก RAM 4 GB, พื้นที่เก็บข้อมูล SSD 40 GB \\ \hline
3 & จอแสดงผล & จอภาพขนาด 16 นิ้วขึ้นไป ความละเอียดระดับ Full HD (1920 $\times$ 1080 พิกเซล) \\ \hline
4 & ระบบเครือข่ายอินเทอร์เน็ต & ความเร็วในการรับส่งข้อมูลไม่น้อยกว่า 100/100 Mbps สำหรับการสกัดข้อมูลข่าวสารและเชื่อมต่อภายนอก \\
\hline
\end{tabular}%
}
\end{table}

\enlargethispage{2\baselineskip}
\subsection{ซอฟต์แวร์}

เครื่องมือด้านซอฟต์แวร์และไลบรารีที่ใช้ในการพัฒนาระบบแสดงดังตารางที่~\ref{tab_software}

\renewcommand{\arraystretch}{0.92}
\begin{longtable}{|C{1.0cm}|L{3.6cm}|C{2.0cm}|L{7.4cm}|}
\caption{ซอฟต์แวร์ที่ใช้ในการพัฒนา}\label{tab_software}\\
\hline
\textbf{ลำดับ} & \textbf{ซอฟต์แวร์ / ไลบรารี} & \textbf{เวอร์ชัน} & \textbf{วัตถุประสงค์การใช้งาน} \\ \hline
\endfirsthead
\caption[]{ซอฟต์แวร์ที่ใช้ในการพัฒนา~\cskkuContinued}\\
\hline
\textbf{ลำดับ} & \textbf{ซอฟต์แวร์ / ไลบรารี} & \textbf{เวอร์ชัน} & \textbf{วัตถุประสงค์การใช้งาน} \\ \hline
\endhead
1 & Ubuntu Linux & 22.04 LTS & ระบบปฏิบัติการหลักของเครื่องแม่ข่ายโพรดักชัน \\ \hline
2 & Python & 3.11.15 & ภาษาโปรแกรมหลักสำหรับการประมวลผลฝั่งเซิร์ฟเวอร์ \\ \hline
3 & uv & 0.6.14 & ตัวจัดการสภาพแวดล้อมเสมือนและแพ็กเกจความเร็วสูง \\ \hline
4 & FastAPI & 0.141.1 & เว็บเฟรมเวิร์กแบบอะซิงโครนัสสำหรับสร้าง REST API \\ \hline
5 & Uvicorn & 0.52.1 & เว็บเซิร์ฟเวอร์ ASGI ประสิทธิภาพสูง \\ \hline
6 & Playwright & 1.62.0 & จำลองเบราว์เซอร์ Chromium ดึงหน้าเว็บ JavaScript และ SPA \\ \hline
7 & Trafilatura & 2.2.0 & สกัดเนื้อหาบทความข่าวและตัดส่วนประกอบที่ไม่เกี่ยวข้อง \\ \hline
8 & PyThaiNLP & 5.3.5 & ตัดคำภาษาไทย กรองคำหยุด และสกัดเอนทิตีชื่อเฉพาะ thainer-v2 \\ \hline
9 & Scikit-learn & 1.6.1 & คำนวณเวกเตอร์ TF-IDF และโมเดลจำแนก LinearSVC \\ \hline
10 & Transformers & 5.15.1 & รันโมเดลภาษา WangchanBERTa สกัดเวกเตอร์ความหมาย 768 มิติ \\ \hline
11 & Tailwind CSS & 3.4.1 & ออกแบบส่วนติดต่อผู้ใช้แบบ Responsive ตามแนวคิด Bento Hub \\ \hline
12 & SQLite3 (WAL Mode) & 3.45+ & ฐานข้อมูลเชิงสัมพันธ์น้ำหนักเบาพร้อมโหมด Write-Ahead Logging \\ \hline
13 & Git & 2.40+ & ระบบควบคุมและจัดการรุ่นของซอร์สโค้ด (Version Control) \\ \hline
14 & Visual Studio Code & 1.90+ & เครื่องมือแก้ไขและพัฒนาโปรแกรมต้นฉบับ \\
\hline
\end{longtable}

% -----------------------------------------------------------
\enlargethispage{2\baselineskip}
\section{แผนการดำเนินงาน}
% -----------------------------------------------------------

แผนการดำเนินงานตลอดระยะเวลาการทำโครงงานตั้งแต่เดือนมิถุนายนถึงเดือนตุลาคม แสดงดังตารางที่~\ref{tab_gantt}

\begin{table}[H]
\centering
\caption{แผนการดำเนินงาน ภาคการศึกษาที่ 2/2568}
\label{tab_gantt}
\renewcommand{\arraystretch}{0.88}
\resizebox{\textwidth}{!}{%
\begin{tabular}{|c|L{9.8cm}|c|c|c|c|c|}
\hline
\textbf{ลำดับ} & \textbf{กิจกรรม / ขั้นตอนการดำเนินงาน} & \textbf{มิ.ย.} & \textbf{ก.ค.} & \textbf{ส.ค.} & \textbf{ก.ย.} & \textbf{ต.ค.} \\ \hline
\multicolumn{7}{|l|}{\textbf{ส่วนที่ 1 การศึกษา ออกแบบ และพัฒนาวิศวกรรมระบบ}} \\ \hline
1 & ศึกษาปัญหา รวบรวมความต้องการ และสืบค้นงานวิจัยที่เกี่ยวข้อง & $\bullet$ & & & & \\ \hline
2 & วิเคราะห์ข้อกำหนดระบบ จัดทำแบบจำลอง และออกแบบสถาปัตยกรรม & $\bullet$ & $\bullet$ & & & \\ \hline
3 & พัฒนาระบบ Scraper Pipeline, Playwright และ Multi-tier Cache & & $\bullet$ & & & \\ \hline
4 & พัฒนาโมเดลจำแนก 9 หมวดหมู่ และระบบระบุประเด็นข่าว 2 ชั้น & & $\bullet$ & $\bullet$ & & \\ \hline
5 & พัฒนาระบบจัดอันดับกระแสข่าว 6 มิติ และแฟ้มสังเคราะห์ไทม์ไลน์ & & & $\bullet$ & & \\ \hline
6 & พัฒนา Bento Hub SPA, WebSocket สด, ระบบค้นหา และ PDPA & & & $\bullet$ & $\bullet$ & \\ \hline
\multicolumn{7}{|l|}{\textbf{ส่วนที่ 2 การทดสอบ ประเมินผล และสรุปรายงานโครงงาน}} \\ \hline
7 & พัฒนาชุดทดสอบอัตโนมัติ Unit Tests และ Integration Smoke Tests & & & & $\bullet$ & \\ \hline
8 & จัดทำแบบทดสอบการยอมรับของผู้ใช้ (UAT) และแบบประเมินความพึงพอใจ & & & & $\bullet$ & \\ \hline
9 & ดำเนินการทดสอบ UAT และเก็บรวบรวมข้อมูลความพึงพอใจจากผู้ใช้จริง & & & & $\bullet$ & $\bullet$ \\ \hline
10 & วิเคราะห์ข้อมูลสถิติ สรุปผลการประเมิน และปรับแต่งระบบขั้นสุดท้าย & & & & & $\bullet$ \\ \hline
11 & จัดทำเอกสารรายงานโครงงานฉบับสมบูรณ์และเตรียมนำเสนอผลงาน & & & & & $\bullet$ \\ \hline
\multicolumn{7}{|l|}{\footnotesize \textit{หมายเหตุ สัญลักษณ์ $\bullet$ หมายถึง ช่วงเวลาดำเนินกิจกรรมในแต่ละขั้นตอนของโครงงาน}} \\ \hline
\end{tabular}%
}
\end{table}

จากตารางที่~\ref{tab_gantt} การดำเนินโครงงานเริ่มต้นจากการศึกษาปัญหาและออกแบบระบบในเดือนมิถุนายนถึงกรกฎาคม ดำเนินการพัฒนาระบบทั้งส่วนหน้าและส่วนหลังในเดือนกรกฎาคมถึงกันยายน และดำเนินการทดสอบระบบ ประเมินผลความพึงพอใจของผู้ใช้ สรุปผล และจัดทำเอกสารรายงานฉบับสมบูรณ์ในเดือนกันยายนถึงตุลาคม

